% Generated draft: requires mathematical and source review.
\documentclass[11pt]{article}
\usepackage[margin=25mm]{geometry}
\usepackage{fontspec}
\setmainfont{Latin Modern Roman}
\usepackage{amsmath,amssymb,mathtools,unicode-math}
\setmathfont{Latin Modern Math}
\usepackage{enumitem,needspace,xurl,hyperref}
\setlength{\parindent}{0pt}
\setlength{\parskip}{5pt}
\setlength{\emergencystretch}{3em}
\newcommand{\N}{\mathbb{N}}
\newcommand{\Z}{\mathbb{Z}}
\newcommand{\Q}{\mathbb{Q}}
\newcommand{\R}{\mathbb{R}}
\newcommand{\C}{\mathbb{C}}
\newcommand{\F}{\mathbb{F}}
\newcommand{\A}{\mathbb{A}}
\newcommand{\PP}{\mathbb{P}}
\newcommand{\E}{\mathbb{E}}
\newcommand{\eps}{\varepsilon}
\newcommand{\dd}{\,\mathrm{d}}
\begin{document}
\section*{UnsolvedMath ID: 801004}
\textbf{Source problem code:} catalogue entry
\subsection{Definitions and mathematical statement}
For a set \(A\) in a Euclidean space, let \(\dim_{\mathrm H}A\) denote its Hausdorff dimension. For each real slope \(m\in\R\), define the scalar orthogonal projection
\[
p_m(x,y)=\frac{x+my}{\sqrt{1+m^2}}.
\]
This identifies orthogonal projection onto the line spanned by \((1,m)\) with a subset of \(\R\), without changing Hausdorff dimension. For the vertical direction, write \(p_\infty(x,y)=y\). For \(S\subseteq\R\cup\{\infty\}\), let \(\Phi(S)\) be the assertion
\[
\exists E\subseteq\R^2\quad
\left(\dim_{\mathrm H}E=1\ \land\
\forall m\in S\;\dim_{\mathrm H}\bigl(p_m(E)\bigr)=0\right).
\]
\textbf{Question.} Can the existence of a Davies example be proved in ZFC? Read literally with real slopes, this asks whether
\[
\mathrm{ZFC}\vdash\Phi(\R).
\]
If the source's phrase ``for all \(\theta\)'' means every line direction, including the vertical direction, the target is instead
\[
\mathrm{ZFC}\vdash\Phi\bigl(\R\cup\{\infty\}\bigr).
\]
The supplied record does not resolve this direction convention. Orthogonal projection is the conventional interpretation used in these formulations; the source does not explicitly specify the projection convention. No measurability, definability, or other regularity condition on \(E\) is imposed in the requested example.

For context, an analytic set is a continuous image of a Polish space. The supplied background attributes to [Mar] (1954) the statement that an analytic \(E\subseteq\R^2\) with \(\dim_{\mathrm H}E=1\) has projections of Hausdorff dimension \(1\) in almost every direction, where ``almost every'' has its usual meaning with respect to Lebesgue measure on direction angles. It attributes to [Dav] (1969) a construction under the continuum hypothesis, \(2^{\aleph_0}=\aleph_1\), of a set of Hausdorff dimension \(1\) whose projections in all stated directions have Hausdorff dimension \(0\). The record phrases this latter claim as CH ``could imply'' such an example; its exact cited theorem and direction convention require source review. These historical assertions are supplied background, not independently verified results in this draft.
\subsection{Short English statement}
Does ZFC prove that there is a set \(E\subseteq\R^2\) of Hausdorff dimension \(1\) whose orthogonal projections all have Hausdorff dimension \(0\)? The source describes an example under CH and asks for existence provable in ZFC. Its use of slopes leaves inclusion of the vertical direction unresolved, and it imposes no regularity condition on \(E\).
\subsection{Sources}
\begin{itemize}
\item Local UnsolvedMath registry: permanent ID \(801004\), exact source problem code \texttt{catalogue entry}, title ``Open Problems in Mathematical Logic — Problem 16.'' The supplied registry labels the problem \texttt{open}; this is a supplied status, not an independently verified literature assessment.
\item Barmpalias, Gao et al., \emph{Open Problems in Mathematical Logic} (2026), Problem 16. Supplied source URL: \url{https://arxiv.org/abs/2608.26628}. The record attributes the problem to Barmpalias, Gao et al., with proposed year 2026. Its provenance records access on 2026-10-16 and extraction by \texttt{arxiv\_latex}; this draft does not claim independent access or verification.
\item [Mar] (1954), as cited in the supplied statement for the analytic-set projection result. No full bibliographic entry or direct reference link is supplied.
\item [Dav] (1969), as cited in the supplied statement for the CH example. No full bibliographic entry or direct reference link is supplied.
\end{itemize}
\end{document}
